\documentclass[final]{beamer}

% ====================
% Packages
% ====================

\usepackage[T1]{fontenc}
\usepackage{lmodern}
\usepackage[orientation=portrait,size=a0,scale=1.0]{beamerposter}
\usetheme{gemini}
\usecolortheme{nott}
\usepackage{graphicx}
\usepackage{booktabs}
\usepackage{tikz}
\usepackage{pgfplots}
\pgfplotsset{compat=1.14}
\usepackage{anyfontsize}
\usepackage{xcolor}
\usepackage[skip=2pt,font=normalsize]{subcaption}
\usepackage{adjustbox}

% ----------------------------------
% For plotting study methodology
% ----------------------------------

\usetikzlibrary{shapes.geometric, arrows}

% Defining Tickz Style
\tikzstyle{startstop} = [rectangle, rounded corners, minimum width=3cm, minimum height=1cm, text centered, text width = 10cm, draw=black, fill=white]

% \tikzstyle{io} = [trapezium, trapezium left angle=70, trapezium right angle=110, minimum width=3cm, minimum height=1cm, text centered, text width = 4.5cm, draw=black, fill=blue!30]

\tikzstyle{process} = [rectangle, minimum width=3cm, minimum height=1cm, text centered, text width = 10cm, draw=black, fill=white]

% \tikzstyle{decision} = [diamond, minimum width=3cm, minimum height=1cm, text centered, draw=black, fill=green!30]

\tikzstyle{arrow} = [ultra thick,->,>=stealth]


% ====================
% Lengths
% ====================

% If you have N columns, choose \sepwidth and \colwidth such that
% (N+1)*\sepwidth + N*\colwidth = \paperwidth
\newlength{\sepwidth}
\newlength{\colwidth}
\setlength{\sepwidth}{0.025\paperwidth}
\setlength{\colwidth}{0.4625\paperwidth}

\newcommand{\separatorcolumn}{\begin{column}{\sepwidth}\end{column}}

% ====================
% Title
% ====================

\title{Neural Combinatorial Optimization for Multi-Agent Path Finding on Graphs}

\author{\textbf{Your Name} \inst{1} \and Advisor/Collaborator \inst{1}}

\institute[shortinst]{\inst{1} Your Institute / Lab}

% ====================
% Footer (optional)
% ====================

\footercontent{%
  \textbf{Keywords:} MAPF \textbullet\ Graph neural networks \textbullet\ Neural combinatorial optimization \hfill
  \textbf{Contact:} \href{mailto:your.email@provider.com}{your.email@provider.com}}
% (can be left out to remove footer)


% ====================
% Logo (optional)
% ====================

% use this to include logos on the left and/or right side of the header:
\logoright{\includegraphics[height=6.3cm]{logos/NHBNI_logo.png}}
\logoleft{\includegraphics[height=7cm]{logos/NLogo-Grey.png}}

% ====================
% Body
% ====================

\begin{document}

\begin{frame}[t]
\begin{columns}[t]
\separatorcolumn

\begin{column}{\colwidth}

% ----------------------------------
% Abstract
% ----------------------------------
  \begin{block}{Abstract}
    Multi-Agent Path Finding (MAPF) on graphs is NP-hard and classical solvers (CBS/EECBS, A*, prioritized planning) trade off optimality and runtime.
    This work asks whether \textbf{graph neural networks} can be trained in a \textbf{(geometric) neural combinatorial optimization} (NCO) setting to improve \textbf{inference-time solution quality} for MAPF.
    We replicate and stress-test a recent learning-to-reweight-edges approach, and identify regimes where learning is unstable or absent.
  \end{block}

% ----------------------------------
% Section: Problem setup
% ----------------------------------
  \begin{block}{Problem setup: Graphs \& MAPF}
    \textbf{Graph:} $G=(V,E,w)$, where $w_e$ is the (learned) edge cost.

    \textbf{MAPF:} agents $\{A_1,\dots,A_k\}$ with start vertices $\{s_i\}\subset V$ and targets $\{t_i\}\subset V$.
    A solution is a set of collision-free paths/policies $\{\pi_i\}$.

    \heading{Neural Combinatorial Optimization (NCO)}
    NCO uses a neural model to parameterize the optimization problem (e.g., edge weights) and relies on a downstream solver to produce a discrete plan.
    \textbf{Geometric NCO} further exploits structure of the solution space to improve the solver output.
  \end{block}

% ----------------------------------
% Section: Research question + objectives
% ----------------------------------
  \begin{alertblock}{Research question \& objectives}
    \textbf{Question:} Can we train GNNs to obtain better MAPF solutions on graphs using a (geometric) NCO approach?

    \vspace{0.4em}
    \begin{itemize}
      \item Learn an edge-cost function $w=\mathrm{NN}_\theta(G,\text{agents})$ that improves a fast heuristic solver (Prioritized Planner, PP).
      \item Benchmark vs. stronger solvers (CBS/EECBS) for sum-of-costs (SOC), total distance, and makespan.
      \item Study failure modes (e.g., instability / ``no learning'') and propose next fixes.
    \end{itemize}
  \end{alertblock}

% ----------------------------------
% Section: Related work
% ----------------------------------
  \begin{block}{Related work}
    \begin{itemize}
      \item \textbf{MAPF survey/benchmarks:} problem definitions and variants \cite{stern2019mapf}.
      \item \textbf{GNN-based MAPF:} temporal-spatial aggregation with attention \cite{zhang2025gatmapf}.
      \item \textbf{Edge-reweighting for MAPF:} learned representations that help a conventional solver \cite{captar2025learned}.
      \item \textbf{Geometric NCO with constraints:} geometric post-processing/updates in discrete spaces \cite{karalias2025geometric}.
    \end{itemize}
  \end{block}

% ----------------------------------
% Section: Method overview
% ----------------------------------
 \begin{block}{Method overview (edge-reweighting + solver-in-the-loop)}
    We learn edge weights $w_e$ so that a fast solver (PP) returns plans closer to an (expensive) reference solver (CBS/EECBS).

    \vspace{0.4em}
    \input{images/figure1}
 \end{block}

% -------------------------------
% Section: Data \& practical challenges
% -------------------------------
\begin{block}{Data \& practical challenges}
  \begin{itemize}
    \item \textbf{Datasets:} synthetic/random graphs; grid maps; real road networks (e.g., SNAP RoadNet-PA; MovingAI).
    \item \textbf{Road networks are huge:} slow to load/visualize; many connected components; high compute for message passing.
    \item \textbf{Edge-level learning:} representation naturally targets edge costs/usage, but supervision is indirect via a planner.
  \end{itemize}
\end{block}

\end{column}

\separatorcolumn

\begin{column}{\colwidth}

% -------------------------------
% Section: Results and discussion
% -------------------------------

\begin{block}{Reference formulation (from \cite{captar2025learned})}
  We follow a solver-in-the-loop objective that learns edge costs.
  Let $y\in\mathbb{N}^{|E|}$ be the edge-usage vector produced by a planner, and $w\in\mathbb{R}^{|E|}$ be learned edge costs.

  \vspace{0.3em}
  \begin{itemize}
    \item Plan cost: $c(w,y)=\sum_{e\in E} w_e\, y_e$.
    \item Discrete solve: $\hat y(w)=\arg\min_{y\in\mathcal{Y}} c(w,y)$ (PP solves this greedily, agent by agent).
    \item Training signal: $L(\hat y)=\frac{1}{|E|}\sum_{e}(\hat y_e-y^*_e)^2$, with $y^*$ the reference plan from CBS/EECBS.
    \item Gradient: black-box differentiation through the planner,
      $\nabla_w L \approx -\frac{1}{\lambda}\big(\hat y(w)-\hat y(w+\lambda\,\partial L/\partial \hat y)\big)$, then backpropagated to $\theta$.
  \end{itemize}
\end{block}

% -------------------------------
% Section: Replication status
% -------------------------------
\begin{block}{Replication status}
  \begin{itemize}
    \item Reproduced the codebase (Flatland railroad variant) released by the authors.
    \item Verified that CBS yields optimal SOC but is expensive; EECBS provides bounded suboptimality with lower runtime.
    \item For small agent counts ($\sim 20$), CBS is often feasible; otherwise EECBS is the practical reference.
  \end{itemize}
\end{block}

% -------------------------------
% Section: Experimental observations
% -------------------------------
\begin{block}{Experimental observations (grids)}
  \heading{Total distance \& makespan}
  Learned edge costs sometimes reduce validation metrics, but curves can be noisy and sensitive to agent count.

  \vspace{0.4em}
  \input{images/figure2}

  \heading{Failure case: ``no learning''}
  In some runs, training stalls or oscillates.
  Observed causes:
  \begin{itemize}
    \item \textbf{Vanishing surrogate gradient:} with a mean-reduced loss the perturbation $\lambda\,\partial L/\partial\hat y$ is $\sim 10^{-3}$, so PP's plan rarely changes and the gradient is zero for most scenarios.
    \item \textbf{Loss--metric mismatch:} the usage loss penalises equally short routes that differ from the reference, so lower loss does not mean lower delay.
    \item \textbf{Low-agent regime:} PP is already close to CBS, leaving little learning signal; validation noise then exceeds the effect.
  \end{itemize}
\end{block}

% -------------------------------
% Section: Conclusions + next steps
% -------------------------------
\begin{alertblock}{Conclusions \& next steps}
  \begin{itemize}
    \item Edge-reweighting is a simple NCO interface to MAPF (continuous parameters, discrete solver).
    \item Replication suggests improvements are possible, but training can be unstable and highly regime-dependent.
    \item Next: (i) focus on makespan-specific objectives, (ii) add geometric post-processing/updates, (iii) visualize solution-space structure on grid maps, (iv) scale via subgraph sampling for road networks.
  \end{itemize}
\end{alertblock}

% -------------------------------
% Section: References
% ------------------------------- 

\begin{block}{References}
  \footnotesize{\bibliographystyle{plain}\bibliography{poster}}
\end{block}

\end{column}
\separatorcolumn



\end{columns}
\end{frame}

\end{document}
